\chapter{Background}\label{chap:background}

This chapter reviews the literature that underpins the thesis. It follows a funnel structure that narrows from general theory to the specific research problem. Section~\ref{sec:bg-mpc} presents the theoretical foundations of \ac{mpc}: the receding-horizon principle, the role of the cost function, constraint handling, and the performance advantages over classical controllers. Section~\ref{sec:bg-apps} examines how \ac{mpc} has been adapted to electromechanical and robotic systems, with emphasis on fast dynamics, multivariable integration, nonlinearity, and computational burden. Section~\ref{sec:bg-wip} then concentrates on \ac{wip} robots, reviewing the state of the art in the control of their underactuated dynamics, disturbance rejection, and stability limits. Section~\ref{sec:bg-gaps} concludes with a critical analysis of the \ac{wip} literature and identifies the gaps that motivate this work.

%======================================================================
\section{Model Predictive Control: Theoretical Foundations}\label{sec:bg-mpc}
%======================================================================

\subsection{The Receding-Horizon Principle and the Cost Function}\label{sec:bg-mpc-mechanics}

\ac{mpc}, also known as receding horizon control, denotes a family of control methods in which an explicit dynamic model of the plant is used to predict future behaviour, and the control action is obtained by solving an optimal control problem online~\cite{rawlings2000,bemporad1999}. The underlying ideas date back to the 1960s, but interest grew rapidly during the 1980s with the industrial algorithms IDCOM and Dynamic Matrix Control, which were conceived for multivariable constrained processes, and with Generalized Predictive Control, which was oriented towards single-variable and adaptive problems~\cite{bemporad1999}. \ac{mpc} has since become an accepted standard for constrained multivariable control in the process industries~\cite{bemporad1999,schwenzer2021}.

The operating mechanism is common to all formulations. At every sampling instant $k$, the state is measured or estimated and an optimisation over a finite prediction horizon is solved; only the first element of the resulting input sequence is applied to the plant, the remainder is discarded, and the procedure is repeated at instant $k+1$ with fresh measurements~\cite{bemporad1999,rawlings2000}. Because the horizon shifts forward with time, it is said to recede. Although each optimisation relies on an open-loop prediction, the repeated re-solution with new measurements provides the feedback characteristic of the controller~\cite{bemporad1999}.

For reference tracking, a representative formulation minimises a quadratic cost over the prediction horizon $N_p$ and the control horizon $N_c$,
\begin{equation}\label{eq:mpc-cost}
J(k)=\sum_{i=1}^{N_p}\left\|\mathbf{y}(k+i\,|\,k)-\mathbf{r}(k+i)\right\|_{\mathbf{Q}}^{2}
+\sum_{i=0}^{N_c-1}\left\|\Delta\mathbf{u}(k+i\,|\,k)\right\|_{\mathbf{R}}^{2},
\end{equation}
subject to the model and constraint set
\begin{equation}\label{eq:mpc-problem}
\begin{array}{ll}
\mathbf{x}(k+i+1\,|\,k)=f\big(\mathbf{x}(k+i\,|\,k),\mathbf{u}(k+i\,|\,k)\big), & \mathbf{u}_{\min}\le\mathbf{u}(k+i\,|\,k)\le\mathbf{u}_{\max},\\[1mm]
\Delta\mathbf{u}_{\min}\le\Delta\mathbf{u}(k+i\,|\,k)\le\Delta\mathbf{u}_{\max}, & \mathbf{x}_{\min}\le\mathbf{x}(k+i\,|\,k)\le\mathbf{x}_{\max},
\end{array}
\end{equation}
where $\mathbf{y}$ is the predicted output, $\mathbf{r}$ the reference, $\mathbf{u}$ the input, and $\mathbf{Q}$ and $\mathbf{R}$ are symmetric positive (semi-)definite weighting matrices. The first term penalises deviation of the predicted outputs from the reference, and the second penalises control effort or its rate of change, so that the weights express an explicit trade-off between tracking accuracy and actuator activity. When the model is linear, the cost is quadratic, and the constraints are linear, the problem is a convex \ac{qp}; an $\ell_1$ or $\ell_\infty$ cost leads instead to a linear program~\cite{bemporad1999,rawlings2000}.

Practical implementations add several elements around this core. When the full state is not measured, a state estimator is required, and non-zero set-points are handled through a steady-state target calculation. Augmenting the model with constant disturbance terms yields integral action and offset-free tracking~\cite{rawlings2000}, a construction that has been adopted in robotic applications~\cite{kamel2017}. A further point relevant to this thesis is the relationship with classical optimal control: for a linear plant without inequality constraints, the infinite-horizon quadratic problem reduces to the \ac{lqr}~\cite{rawlings2000}. \ac{mpc} can therefore be regarded as a constrained, finite-horizon, receding generalisation of \ac{lqr}, which makes \ac{lqr} the natural benchmark for any \ac{mpc} design.

The prediction model is the essential component of the controller~\cite{rawlings2000}. In contrast to conventional control design, \ac{mpc} shifts the engineering effort from selecting a controller structure towards modelling the plant, and because the tuning parameters retain a physical interpretation, tuning is comparatively intuitive~\cite{schwenzer2021}. State-space models are the usual choice in \ac{mpc} theory, since they extend directly to multivariable systems and are convenient for closed-loop analysis and online computation~\cite{rawlings2000}. Nonlinear prediction models can improve forecasts when the plant operates over a large region of the state space, but they transform the optimisation into a nonconvex problem with more decision variables~\cite{rawlings2000,kuhne2004}.

\subsection{Constraint Handling}\label{sec:bg-mpc-constraints}

The feature that most clearly distinguishes \ac{mpc} from unconstrained optimal control is that limits on inputs, input rates, states, and outputs enter the optimisation directly as constraints~\cite{rawlings2000,bemporad1999}. Constrained optimal control was not new when \ac{mpc} emerged; the contribution lay in the moving-horizon implementation, which re-solves the constrained problem at every sample and thereby yields a feedback law~\cite{rawlings2000}. The controller thus plans with explicit knowledge of actuator saturation and admissible operating regions, instead of computing an unconstrained command that is subsequently clipped. In mobile robotics, this property has been cited as a primary motivation for \ac{mpc}: control actions that respect actuator limits are generated automatically, and state constraints can confine the robot to a safe region of operation~\cite{kuhne2005}.

Hard state or output constraints may render the optimisation infeasible, for instance after a large disturbance. Soft constraints, in which slack variables are penalised in the cost (an exact-penalty formulation), preserve feasibility while still discouraging violations~\cite{rawlings2000,schwenzer2021}. Theoretical guarantees such as recursive feasibility and closed-loop stability have been established for linear formulations~\cite{schwenzer2021} and exploited, for example, in mobile manipulation~\cite{buizza2018}.

Because plant--model mismatch and disturbances are unavoidable, a substantial body of work addresses robustness. Bemporad and Morari survey the uncertainty descriptions used in the \ac{mpc} literature (including bounded disturbances and parametric uncertainty) together with techniques for robust constraint satisfaction, stability, and performance, and they emphasise closed-loop prediction, i.e.\ accounting for future feedback in the predictions in order to limit conservatism~\cite{bemporad1999}. Robust formulations based on worst-case cost minimisation are revisited in Section~\ref{sec:bg-apps-nonlinear}.

\subsection{Performance Advantages over Classical Controllers}\label{sec:bg-mpc-advantages}

Classical feedback controllers such as \ac{pid} act on the present and past error and are therefore reactive, whereas \ac{mpc} uses a model to anticipate deviations from the reference~\cite{schwenzer2021}. Although \ac{pid} control remains the most widespread industrial solution, finding a suitable parametrisation is often difficult for nonlinear or time-variant systems~\cite{schwenzer2021}. \ac{mpc}, by contrast, is naturally multivariable: coupled inputs and outputs are treated in a single optimisation, without the need for separate loops and decoupling compensators~\cite{rawlings2000,bemporad1999}. Its tuning parameters are directly related to an explicit performance criterion, and constraints are accommodated in a straightforward manner~\cite{kuhne2005}.

In some cases, \ac{mpc} can stabilise systems that smooth, time-invariant feedback cannot. Nonholonomic wheeled mobile robots are the standard example: by Brockett's condition, no continuously differentiable time-invariant feedback law can stabilise them, whereas \ac{mpc} naturally produces a discontinuous control law~\cite{kuhne2004}. Schwenzer et al.\ accordingly note that \ac{mpc} can control systems that conventional feedback controllers cannot~\cite{schwenzer2021}.

The price of these advantages is online computation. In the process industries, where \ac{mpc} originated, sampling times were once measured in hours, and later in seconds or minutes~\cite{schwenzer2021,wang2010}. This is the central obstacle when the technique is transferred to electromechanical plants, and it motivates the next section.

%======================================================================
\section{MPC for Electromechanical and Robotic Systems}\label{sec:bg-apps}
%======================================================================

Applying \ac{mpc} to physical electromechanical systems stresses the assumptions that make it successful in slow processes. The literature reviewed below, covering mobile robots, manipulators, aerial vehicles, and road vehicles, reveals four recurring challenges: fast dynamics and computational burden, multivariable and constraint coupling, nonlinearity, and model uncertainty. Interest in the technique is nevertheless growing quickly; in legged robotics, for example, 19 of the 28 workshop papers presented at ICRA~2022 referred to model predictive or receding-horizon methods~\cite{katayama2023}.

\subsection{Fast Dynamics and Computational Burden}\label{sec:bg-apps-fast}

A widely recognised limitation of \ac{mpc} is that it is usually confined to applications with slow dynamics, because an optimisation problem must be solved within each sampling interval~\cite{wang2010}. K\"uhne et al.\ make this point for wheeled mobile robots: \ac{mpc} is accepted for slow, constrained, multivariable plants, but fast or nonlinear dynamics limit its applicability because of the required online computation~\cite{kuhne2004}. In a subsequent study on trajectory tracking, the same group presents the computational burden as the principal drawback of \ac{mpc}, argues that faster processors and efficient algorithms make robotic applications feasible, and analyses real-time viability explicitly~\cite{kuhne2005}. Buizza Avanzini et al.\ reach the same conclusion for mobile manipulators: the cost of \ac{mpc} was traditionally incompatible with the high sampling rates of robots, but tailored solution algorithms now permit online control of such fast systems~\cite{buizza2018}.

Several strategies for reducing the burden have been reported. Wang and Boyd exploit the structure of the \ac{qp} that arises in \ac{mpc}, so that the cost per interior-point step grows linearly rather than cubically with the horizon length. Their custom solver is roughly two orders of magnitude faster than a generic optimiser; a problem with 12 states, 3 inputs, and a horizon of 30 steps is solved in about 5~ms, which permits operation at 200~Hz~\cite{wang2010}. The same work notes that precomputing the control law offline (explicit \ac{mpc}) and evaluating it as a lookup table is practical only for small problems, with state and input dimensions of about five or less, few constraints, and short horizons~\cite{wang2010,schwenzer2021}. Schwenzer et al.\ identify the computational burden as a persistent design factor despite continued progress~\cite{schwenzer2021}. In legged robotics, high-dimensional whole-body models lead to large nonlinear programs with costly dynamics derivatives, and the response has been to adopt reduced-order models and Riccati-recursion-based Newton-type solvers~\cite{katayama2023}. Architectural measures are also used: in the \ac{uav} framework of Kamel et al., a reliable low-level controller runs as an inner loop on microcontroller-based flight hardware, whereas the \ac{mpc} outer loop runs on a more powerful but less reliable onboard computer, so that safety-critical functions survive a failure of the high-level computer~\cite{kamel2017}.

\subsection{Multivariable Integration and Constraint Coupling}\label{sec:bg-apps-multivariable}

Robotic platforms generally combine several actuated degrees of freedom whose motion is subject to coupled physical limits. Buizza Avanzini et al.\ apply constrained \ac{mpc} to a mobile manipulator, treating base and arm as a single redundant system. Acceleration, velocity, and position limits, collision avoidance, and field-of-view constraints are all extended over the prediction horizon while the problem remains linear-quadratic, and offset-free tracking with convergence and recursive feasibility guarantees is obtained and verified experimentally~\cite{buizza2018}. They also point out that instantaneous (single-step) constrained optimisation is blind to future system evolution, so that a feasible solution at one instant may lead to an infeasible configuration later; the prediction horizon of \ac{mpc} removes this weakness~\cite{buizza2018}. A comparable breadth is reported for autonomous ground vehicles, for which \ac{mpc} has been used for tracking, planning, and the coordination of multiple vehicles under centralised and distributed architectures~\cite{yu2021}. For aerial robots, linear and robust \ac{mpc} with a disturbance model for offset-free tracking has been applied to multi-rotor vehicles, and nonlinear \ac{mpc} to the lateral-directional control of fixed-wing aircraft, with integration into the Robot Operating System and experimental validation~\cite{kamel2017}.

\subsection{Nonlinearity, Model Uncertainty, and Model Fidelity}\label{sec:bg-apps-nonlinear}

The accuracy of the prediction model directly bounds the quality of the control, which creates tension between fidelity and computational cost. Several compromises have been proposed. K\"uhne et al.\ avoid solving a nonconvex nonlinear program at each step by successively linearising an error model of the mobile robot, obtaining a linear time-varying description that results in a \ac{qp}~\cite{kuhne2004}, and they compare linear and nonlinear \ac{mpc} for trajectory tracking in simulation~\cite{kuhne2005}. Cheng et al.\ treat the time-varying forward speed of an autonomous vehicle by building a \ac{lpv} lateral model over a polytope of speeds. Their robust \ac{mpc} minimises the worst-case infinite-horizon quadratic cost through a set of \ac{lmi}, and it is validated in co-simulation and on a real bus~\cite{cheng2021}. Rokonuzzaman et al.\ observe that kinematic vehicle models become inadequate at higher speeds because of unmodelled dynamics, whereas high-fidelity dynamic models incur a high computational cost. They propose a neural-network model of the vehicle dynamics inside the \ac{mpc}, which outperforms \ac{mpc} with conventional models in their evaluation and is compared against linear and nonlinear \ac{mpc}~\cite{rokonuzzaman2021}. When analytical models are unavailable, system identification has been used to obtain the prediction model, which is then combined with robust linear \ac{mpc}~\cite{kamel2017}. In legged robotics, the choice between reduced-order and whole-body models is regarded as the central design decision, and frequent replanning from the current state is reported to improve robustness against disturbances~\cite{katayama2023}.

\subsection{Synthesis}

Taken together, these studies show that applying \ac{mpc} to an electromechanical system requires balancing model fidelity, constraint richness, and solution time within a sampling interval dictated by the plant dynamics. The reported remedies (successive linearisation, reduced-order or learned models, structure-exploiting solvers, explicit laws, and cascaded architectures) are largely application-specific. In the studies reviewed above, the validated applications also involve plants of comparatively moderate bandwidth, such as vehicles, manipulators, and quadrotors with a dedicated inner loop. A \ac{wip} robot is more demanding in that its open-loop dynamics are unstable and underactuated, its actuators saturate, and its admissible tilt range is small, so that stabilisation demands a high sampling rate. These properties combine every difficulty identified above, as discussed in the next section.

%======================================================================
\section{MPC for Wheeled Inverted Pendulum Robots}\label{sec:bg-wip}
%======================================================================

\subsection{Dynamics: Underactuation, Instability, and Nonholonomy}\label{sec:bg-wip-dynamics}

A \ac{wip} robot is driven by two independent wheel motors that carry an inverted-pendulum-like body. Its control problem has two coupled objectives: trajectory tracking in the plane, and stabilisation of the unstable body posture, the latter being a prerequisite for the former~\cite{yue2018tsmc}. Only the wheel torques are available, so the body pitch is not directly actuated, which makes the system underactuated; at the same time, the rolling-without-slipping condition imposes a nonholonomic constraint, so that by Brockett's theorem the robot cannot be stabilised to an arbitrary configuration by continuous time-invariant state feedback~\cite{yue2018ijcas}. Unlike many other underactuated mechanical systems, the equilibrium of the \ac{wip} (the upright posture) is unstable. Because the drive motors are limited in power and torque, the tilt angle must be kept within a small range: once it is exceeded, the saturated actuators can no longer restore the upright posture~\cite{yue2018tsmc}.

Models are commonly derived from first principles (for example, with the Newton--Euler approach) and linearised about the upright equilibrium. A recurring simplification is the decomposition of the dynamics through the sum and difference of the wheel torques into a ``pendulum'' subsystem (tilt, tilt rate, position, and velocity, driven by the torque sum) and a ``rotation'' subsystem (heading and heading rate, driven by the torque difference)~\cite{minouchehr2015,yue2018tsmc}. The underactuated pendulum subsystem is the critical one, since its single input must simultaneously regulate forward velocity and body stability~\cite{yue2018tsmc}. In wheel-legged (wheeled-biped) robots, the leg configuration changes the height and distribution of the \ac{com} and the moments of inertia, so the linearised model depends on posture. The wheeled-biped Ascento reduces to a \ac{wip} at a fixed leg geometry~\cite{klemm2019}, whereas the SR600-II robot represents the dynamics as a \ac{wip} model whose matrices depend on the upper-body pose~\cite{zhang2022}.

\subsection{Baseline Controllers: LQR and Gain-Scheduled Variants}\label{sec:bg-wip-lqr}

\ac{lqr} is the principal baseline in the wheeled-biped work reviewed here. In the Ascento robot, the stabilising controller is an \ac{lqr} acting on the tilt angle, tilt rate, velocity, and yaw rate, fed by a Kalman filter that fuses inertial and encoder measurements; the gains are interpolated across leg configurations to relax the fixed-geometry assumption, and the robot withstands strong impulsive hits as well as sustained pushes, pulls, and added loads in experiments~\cite{klemm2019}. The authors state that the stability guarantee of this setup holds only as long as noise, disturbances, modelling errors, and actuator saturation do not influence the dynamics, and they explicitly identify \ac{mpc} as an interesting direction for further control development on the platform~\cite{klemm2019}.

Zhang et al.\ extend this approach to the SR600-II wheel-biped transformable robot with a parameter-varying \ac{lqr} whose matrices are recomputed as the upper-body pose changes, running with a 5~ms control step. In simulation, the controller recovers from an impact in about 4~s and tolerates a maximum impact force of 132~N, whereas the \ac{pid} baseline falls above 80~N; hardware experiments demonstrate driving while squatting and extending, and wheel-to-foot transformation~\cite{zhang2022}. Safety limits in that work are expressed as geometric balance conditions on the \ac{com} projection and tilt angle, and they are enforced by the mode-switching logic (the \ac{lqr} is turned off during the critical transition) rather than inside the controller's optimisation~\cite{zhang2022}. These studies illustrate both the strength and the limitation of \ac{lqr}: it is simple, fast, and hardware-proven, but it is optimal only for the unconstrained linear problem and therefore cannot encode actuator saturation or tilt limits.

\subsection{MPC-Based Approaches}\label{sec:bg-wip-mpc}

Compared with the \ac{lqr} work above, \ac{mpc} for \ac{wip} robots is sparsely studied; Yue et al.\ observe that, although \ac{mpc} has been applied to many mechanical systems, only a few studies address underactuated plants~\cite{yue2018tsmc}. Three studies in the reviewed set apply \ac{mpc} directly to \ac{wip}-type vehicles, and their characteristics are compared in Table~\ref{tab:wip-summary}.

Minouchehr and Hosseini-Sani design \ac{mpc} for a two-wheeled inverted pendulum by linearising the Newton--Euler model about the upright equilibrium and splitting it with a decoupling unit into the pendulum and rotation subsystems, each controlled by an independent predictive controller. The optimisation is solved in closed form as a least-squares problem, without constraints. In simulation on the nonlinear model, with a 0.1~s sampling time, sinusoidal and step-like external forces are applied, and the \ac{mpc} yields lower mean-square error and lower control effort than an \ac{lqr} with the same weighting matrices. The authors attribute the improvement to the capacity of \ac{mpc} to compensate measurable disturbances~\cite{minouchehr2015}. It should be noted that an unconstrained, linear \ac{mpc} with the same weights is closely related to \ac{lqr} (Section~\ref{sec:bg-mpc-mechanics}), so this comparison does not isolate the principal advantage of \ac{mpc}, namely constraint handling.

Yue et al.\ propose a hierarchical scheme for a \ac{wip} vehicle that also lacks velocity measurements. An \ac{eso} estimates the velocities and the lumped uncertainty from position measurements alone, with uniform ultimate boundedness established through a Lyapunov argument. \ac{mpc} is applied to the underactuated longitudinal subsystem to track the longitudinal velocity while keeping the body upright, with explicit constraints on velocity, tilt angle and rate, torque, and torque increment, formulated as a \ac{qp}. An adaptive robust controller handles the actuated rotational subsystem, and a saturated trajectory generator accommodates velocity limits and the nonholonomic constraint. Simulation with sinusoidal and pulse disturbances shows robust tracking of a trifolium path (sampling time 0.1~s, $N_p=20$, $N_c=6$)~\cite{yue2018tsmc}.

In a second study, Yue et al.\ apply approximate input--output feedback linearisation before the \ac{mpc} in order to lighten the computation and improve real-time performance, again combined with a saturated trajectory generator that guarantees convergence of the posture errors to compact sets. Bounds on torque, torque increment, tilt angle and rate ($\pm 0.1$~rad and $\pm 0.1$~rad/s in the simulation), and velocity are enforced. Robustness is assessed in simulation against initial errors, a large pulse disturbance, and random noise on a circular path~\cite{yue2018ijcas}.

\begin{table}[ht]
\centering
\footnotesize
\setlength{\tabcolsep}{4pt}
\renewcommand{\arraystretch}{1.2}
\begin{tabular}{@{}>{\raggedright\arraybackslash}p{1.9cm}>{\raggedright\arraybackslash}p{2.4cm}>{\raggedright\arraybackslash}p{2.9cm}>{\raggedright\arraybackslash}p{2.4cm}>{\raggedright\arraybackslash}p{2.3cm}>{\raggedright\arraybackslash}p{1.8cm}@{}}
\toprule
\textbf{Study} & \textbf{Platform / model} & \textbf{Controller} & \textbf{Constraints in the controller} & \textbf{Disturbance treatment} & \textbf{Validation} \\
\midrule
Minouchehr and Hosseini-Sani~\cite{minouchehr2015} & Two-wheeled inverted pendulum, linearised and decoupled & Two linear MPCs; compared with LQR & None (closed-form solution) & Measurable external forces (sinusoidal, step) & Simulation, $T_s=0.1$~s \\
Yue et al.~\cite{yue2018tsmc} & WIP vehicle, nonlinear reduced-order model & MPC (longitudinal) with adaptive robust control (rotation) and ESO & Velocity, tilt and rate, torque, torque increment & ESO-estimated lumped disturbance; sinusoidal and pulse & Simulation, $T_s=0.1$~s \\
Yue et al.~\cite{yue2018ijcas} & WIP vehicle, nonlinear model & Feedback linearisation with MPC and saturated trajectory generator & Torque, torque increment, tilt and rate, velocity & Pulse disturbance and random noise & Simulation, $T_s=0.1$~s \\
Klemm et al.~\cite{klemm2019} & Ascento wheeled biped & LQR with Kalman filter; interpolated gains & None (stability guarantee excludes saturation) & Hits, pushes, added loads & Hardware \\
Zhang et al.~\cite{zhang2022} & SR600-II wheel-biped & Parameter-varying LQR, 5~ms step & Geometric balance limits via mode logic & Impact forces (simulation) & Simulation and hardware \\
\bottomrule
\end{tabular}
\caption{Summary of the reviewed control studies on WIP-type platforms. $T_s$: sampling time.}
\label{tab:wip-summary}
\end{table}

\subsection{Disturbance Rejection and Stability Limits}\label{sec:bg-wip-dist}

Disturbance rejection is treated in the \ac{mpc}-based \ac{wip} literature in three ways. In the first, measurable external forces are compensated by the predictive law itself~\cite{minouchehr2015}. In the second, an observer estimates the unmeasured velocities and the lumped disturbance, which are then fed to the controller~\cite{yue2018tsmc}. In the third, feedback regulation alone recovers the vehicle after a pulse or random perturbation~\cite{yue2018ijcas}. In each case the evidence is numerical, whereas disturbance rejection for the \ac{lqr}-based robots has been demonstrated physically, through hits with an object and sustained loads on Ascento~\cite{klemm2019}, and through quantified impact limits for the SR600-II~\cite{zhang2022}.

The stability limits of the \ac{wip} are governed by the interaction of the tilt range with actuator saturation. The reviewed \ac{mpc} studies encode this interaction as explicit bounds on tilt angle, tilt rate, and torque inside the optimisation, so that the predicted trajectory remains in the region from which the saturated actuators can still recover~\cite{yue2018tsmc,yue2018ijcas}. By contrast, \ac{lqr} designs rely on tuning and on the nominal guarantee~\cite{klemm2019}, or on geometric conditions enforced externally~\cite{zhang2022}. No reviewed work, however, quantifies how much of the recoverable disturbance envelope (for example, the maximum recoverable impulse or the region of attraction) is gained by the constrained formulation.

%======================================================================
\section{Critical Analysis and Research Gaps}\label{sec:bg-gaps}
%======================================================================

The foregoing review establishes that \ac{mpc} offers a principled framework for the \ac{wip} problem, since its core strengths (explicit treatment of actuator and tilt constraints, prediction, and multivariable handling) correspond to the central difficulties of the platform, and that computational advances have made the technique plausible for fast robotic systems. Nevertheless, a comparison of the \ac{wip} studies in Table~\ref{tab:wip-summary} reveals several limitations.

\begin{enumerate}
\item \textbf{Absence of experimental validation.} All three \ac{mpc} studies on \ac{wip} platforms are simulation-only~\cite{minouchehr2015,yue2018tsmc,yue2018ijcas}, whereas the \ac{lqr}-based wheeled-biped controllers are validated on hardware~\cite{klemm2019,zhang2022}. Effects that dominate real implementations, such as sensor noise, actuator dynamics, delays, and model mismatch, are therefore not examined for \ac{mpc}.
\item \textbf{Unreported computational feasibility.} All three \ac{mpc} studies use a sampling time of 0.1~s, which is twenty times slower than the 5~ms control step reported for a hardware \ac{lqr} on a wheel-biped robot~\cite{zhang2022}. None of them reports an execution time or the processor used. Feedback linearisation is introduced explicitly to reduce the burden~\cite{yue2018ijcas}, but its effect is not quantified. Whether the \ac{qp} can be solved within the short sampling interval demanded by an unstable pendulum, on resource-constrained embedded hardware, remains open, even though the literature shows that structure-exploiting solvers and cascaded architectures make this plausible for problems of comparable size~\cite{wang2010,kamel2017}.
\item \textbf{Inconclusive benchmarking against \ac{lqr}.} The only direct \ac{mpc}--\ac{lqr} comparison is a single simulation study in which \ac{mpc} is unconstrained and uses the same weights as the \ac{lqr}~\cite{minouchehr2015}. The metric is mean-square error, and no analysis addresses recovery time, maximum recoverable disturbance, control effort under saturation, or robustness to parameter variation. The practical benefit of \ac{mpc} over a well-tuned \ac{lqr} on the same platform is thus unproven.
\item \textbf{Limited characterisation of disturbance rejection and stability limits.} Disturbance handling relies either on measurable disturbances or on observer-based estimates, and it is evaluated on a small set of disturbance signals without a systematic study of magnitude, duration, or direction. The stability margin gained from constraint-aware prediction is not quantified (Section~\ref{sec:bg-wip-dist}).
\item \textbf{Neglect of configuration-varying wheel-legged platforms.} Wheeled-biped robots are controlled by \ac{lqr} with interpolated or parameter-varying gains~\cite{klemm2019,zhang2022}, and the authors of Ascento explicitly propose \ac{mpc} as future work~\cite{klemm2019}. Meanwhile, \ac{mpc} is becoming standard in legged robotics~\cite{katayama2023}, and successive linearisation offers a way to handle configuration-dependent models at the cost of a \ac{qp}~\cite{kuhne2004}. No study applies \ac{mpc} to this class of \ac{wip}-based platforms.
\end{enumerate}

In summary, the literature demonstrates the theoretical suitability of \ac{mpc} for \ac{wip} robots but not its practical viability. An experimentally validated \ac{mpc} stabilisation and tracking controller for a \ac{wip}-based robot, which is solved in real time on embedded hardware, explicitly exploits actuator and tilt constraints, and is benchmarked against an \ac{lqr} baseline under identical conditions with quantified disturbance rejection, has not been reported. This gap defines the objective of the present thesis.
% TODO: tailor the final sentence above to the exact objective and contributions of the thesis.
